\chapter{Introduction}\label{ch:introduction}

\section{Motivation}

Fluid dynamics governs a vast range of physical phenomena, from atmospheric circulation and ocean currents to industrial flow systems and aerodynamic design. At the heart of many of these processes lies the Navier-Stokes system of partial differential equations, which describes the conservation of mass and momentum in a fluid medium. Despite their well-established theoretical foundation, the Navier-Stokes equations remain notoriously difficult to solve in practice. Their nonlinear, coupled structure gives rise to complex multi-scale behaviour that often requires high-resolution numerical methods and substantial computational resources. As a result, large-scale simulation campaigns can become prohibitively expensive, particularly when many trajectories must be evaluated repeatedly for design, optimisation, or scientific analysis.

Recent advances in machine learning have created new opportunities to accelerate fluid simulation through the use of \emph{surrogate models}. Rather than solving the governing equations directly, a surrogate learns to predict future system states from previously observed states, transfering the cost of simulation into a training phase and enabling orders-of-magnitude faster inference. Such models have demonstrated promising results on individual fluid dynamics benchmarks and increasingly complex flow regimes.

However, many existing surrogate models are developed and evaluated within relatively narrow settings, often focusing on a single dataset, flow configuration, or physical regime. Real-world fluid dynamics encompasses a much broader spectrum of behaviours, including laminar and turbulent flows, compressible and incompressible dynamics, diverse boundary conditions. Understanding whether a single neural architecture can learn effectively across this heterogeneous landscape is therefore an important step toward more generally applicable fluid dynamics surrogate model.

\section{Problem Statement}

Building a foundation model for Navier-Stokes dynamics presents several fundamental challenges. First, the governing system exhibits qualitatively different behaviour depending on the flow regime. Laminar and turbulent flows, compressible and incompressible conditions, forced and unforced dynamics, and diverse boundary geometries each impose distinct structural constraints on the solution manifold. A model trained on a narrow subset of these conditions will generalise poorly to others, making breadth of training data a critical prerequisite.

Second, turbulent fluid dynamics are intrinsically stochastic in character at the macroscopic level. Even though the Navier-Stokes equations are deterministic, extreme sensitivity to initial conditions means that small errors in the initial state can rapidly propagate into qualitatively different trajectories. A deterministic surrogate model, constrained to produce a single point estimate of the future state, may therefore be fundamentally ill-suited to capture this distributional uncertainty. A probabilistic formulation that learns a conditional distribution over plausible futures is a natural alternative, but imposes its own challenges in terms of architecture design, training stability, and sampling efficiency.

Third, accurate long-horizon prediction requires not only low single-step error but also that predictions remain physically consistent over many autoregressive steps. Error accumulation during rollout is a well-documented failure mode of data-driven surrogate models, and mitigating it demands both architectural choices that promote temporal stability and training procedures that are robust to distributional shift between teacher-forced training and free-running inference.

\section{Approach}

This work addresses these challenges through a comparative study of two surrogate modelling approaches for two-dimensional Navier-Stokes dynamics, trained jointly on a heterogeneous multi-source dataset spanning nine distinct flow configurations.

Both approaches share a common hierarchical backbone, combining a Swin Transformer \citep{Liu2021SwinTA} with a factorised space-then-time attention mechanism embedded in a UNet-style encoder-decoder architecture. This design enables efficient multi-scale feature extraction across spatial and temporal dimensions, with skip connections preserving fine-grained information across resolution levels. The backbone is adapted from the standard image-classification Swin Transformer to handle sequential, spatiotemporal input and to return predictions at the original spatial resolution.

The two modelling strategies differ in their prediction formulation. The \textbf{deterministic model} directly maps a block of five historical velocity field states to the subsequent five states, minimising a mean-squared error objective. The \textbf{probabilistic model} adopts a flow-matching formulation \citep{Lipman2022FlowMatching}, learning to transport samples from a structured Gaussian prior to the target distribution of future states conditioned on the same historical input. By sampling from the prior at inference time, the probabilistic model produces a distribution of plausible futures rather than a single prediction, offering a principled way to represent the inherent uncertainty of turbulent dynamics.

The pre-training dataset is assembled from nine datatsets, covering compressible and incompressible flows, laminar and turbulent regimes, periodic and Dirichlet boundary conditions, and a wide range of initialisation schemes. Together, the datasets contribute approximately 21 TB of simulated two-dimensional velocity field trajectories.

\section{Contributions}

The main contributions of this work are as follows:

\begin{enumerate}
    \item \textbf{A custom spatiotemporal backbone for fluid dynamics.} We design and implement a hierarchical Swin Transformer backbone extended with factorised space-then-time attention blocks and a UNet-style encoder-decoder structure, specifically adapted for multi-dataset pretraining on two-dimensional Navier-Stokes velocity fields.

    \item \textbf{A probabilistic flow-matching surrogate model.} We develop and train a conditional flow-matching model that learns a distribution over future fluid states, conditioned on historical velocity inputs. The model employs a structured spatially-correlated Gaussian prior and integrates over flow-time using a fixed-step Euler scheme, and is shown to outperform deterministic baselines on the large majority of evaluated datasets.

    \item \textbf{Large-scale multi-source pretraining.} We curate, preprocess, and unify nine heterogeneous datasets into a common training pipeline, covering a broad spectrum of Navier-Stokes regimes and providing a foundation for generalisation across flow conditions.

    \item \textbf{Systematic evaluation across accuracy, stability, and spectral fidelity.} We evaluate both model families under single-step, autoregressive rollout, and spectral diagnostics settings, providing a structured comparison of deterministic and probabilistic modelling strategies for fluid dynamics.
\end{enumerate}

\section{Report Structure}

The remainder of this report is organised as follows. Chapter~\ref{ch:problem} formally defines the Navier-Stokes system, the machine learning objective, and the dual modelling framework. Chapter~\ref{ch:datasets} describes the nine pretraining datasets and the preprocessing pipeline. Chapter~\ref{ch:models} details the shared backbone architecture, the mathematical formulation of the attention mechanisms, and the specific design of the deterministic and flow-matching model variants. Chapter~\ref{ch:training} describes the training framework, hardware platform, optimisation strategy, and validation methodology. Chapter~\ref{ch:results} presents and discusses the experimental results across single-step prediction, autoregressive rollout, and spectral analysis.